\chapter{Metodologia}
\label{chap:methodology}

\paragraph{Riassunto del capitolo.} Il capitolo formalizza la metodologia
condivisa fra i tre task: feature engineering deterministico (9 feature
per la regressione, 10 per la classificazione), \texttt{StandardScaler}
fittato solo sul training, strategia di cross-validation
\emph{coerente con il task} (\acrshort{logo} per Task 1 e Task 3,
\texttt{StratifiedGroupKFold} per Task 2) per evitare leakage,
metriche standard ($R^2$/MSE/MAE per la regressione, Accuracy/F1/AUC
per la classificazione) e tuning di iperparametri tramite GridSearchCV.
Il filo logico è: feature $\to$ scaling $\to$ split coerente $\to$
benchmark $\to$ tuning $\to$ valutazione single-shot sul fold di test.

I tre task condividono la stessa \emph{ossatura}: feature engineering
deterministico, scaling, training di un benchmark di modelli, valutazione
con una strategia di cross-validation \emph{coerente con il task}, e
hyperparameter tuning condizionale alla scelta del modello vincente. Il
capitolo formalizza questi cinque passi, evidenziando i punti dove la
metodologia diverge fra regressione (Task 1, 3) e classificazione (Task 2).

\section{Feature engineering}

\subsection{Task di regressione (1, 3)}

Le features sono nove, definite in
\codefile{src/data/features.py::build\_regression\_features}:

\begin{table}[H]
  \centering
  \caption{Feature per Task 1 e Task 3.}
  \label{tab:regression-features}
  \begin{tabularx}{\linewidth}{lX}
    \toprule
    \textbf{Feature} & \textbf{Definizione e motivazione} \\
    \midrule
    \texttt{Aging} & livello discreto $\{0..4\}$ \\
    \texttt{Temperature} & temperatura in \si{\celsius} \\
    \texttt{SOC} & livello discreto $\{0..4\}$ \\
    \texttt{Frequency} & frequenza in \si{\hertz} \\
    \texttt{log\_Freq} & $\log_{10}(\texttt{Frequency})$, dato che la
                        griglia è log-spaziata e i fenomeni elettrochimici
                        sono separati su decadi \\
    \texttt{inv\_Temp} & $1/(T + 273.15)$, forma Arrhenius della
                        dipendenza termica della cinetica \\
    \texttt{Aging\_x\_Temp} & interazione: l'effetto della temperatura
                              sull'impedenza dipende dallo stato di
                              invecchiamento \\
    \texttt{SOC\_x\_Temp}   & interazione: la diffusione dipende
                              dall'occupazione degli stati di intercalazione \\
    \texttt{SOC\_x\_Aging}  & interazione: cellule più vecchie mostrano una
                              pendenza diversa al variare di SOC \\
    \bottomrule
  \end{tabularx}
\end{table}

I target sono \emph{multi-output}: $y = (\mathrm{Re}(Z), \mathrm{Im}(Z))$.
Anche i target sono standardizzati con \texttt{StandardScaler} fittato sul
solo training; le metriche finali sono calcolate sulla scala originale
tramite \texttt{inverse\_transform}.

\subsection{Task di classificazione (2)}

Per il Task 2 l'\emph{impedenza diventa un input}: ciò che era target nei
task di regressione qui descrive la cella, e l'output è la classe
\emph{Young}/\emph{Old}. Si usano dieci features:

\begin{table}[H]
  \centering
  \caption{Feature per Task 2.}
  \label{tab:classification-features}
  \begin{tabularx}{\linewidth}{lX}
    \toprule
    \textbf{Feature} & \textbf{Definizione} \\
    \midrule
    \texttt{Temperature} & \si{\celsius} \\
    \texttt{Frequency} & \si{\hertz} \\
    \texttt{log\_Freq} & $\log_{10}(\texttt{Frequency})$ \\
    \texttt{Z\_real} & $\mathrm{Re}(Z)$ in \milliohm{} \\
    \texttt{Z\_imag} & $\mathrm{Im}(Z)$ in \milliohm{} \\
    \texttt{Z\_magnitude} & $\sqrt{Z_\text{real}^{2} + Z_\text{imag}^{2}}$ \\
    \texttt{Z\_phase} & $\arctan2(Z_\text{imag}, Z_\text{real})$ \\
    \texttt{inv\_Temp} & $1/(T + 273.15)$ \\
    \texttt{Z\_real\_x\_Temp} & interazione resistiva $\times$ termica \\
    \texttt{sqrt\_Freq} & $\sqrt{\texttt{Frequency}}$, vicino al
                          comportamento Warburg ($Z \propto 1/\sqrt{\omega}$) \\
    \bottomrule
  \end{tabularx}
\end{table}

Il target è
$\texttt{Age\_class} = \mathbb{1}\!\big(\texttt{Aging} > \texttt{young\_max\_aging}\big)$,
con \texttt{young\_max\_aging} letto da \codefile{config/config.yaml}
(default $= 2$). \texttt{Aging} è quindi escluso dalle feature: usarlo
come input renderebbe il problema banale.

\section{Scaling}

In tutti i task usiamo \codefile{StandardScaler} fittato \emph{solo sul
training}. Per i task di regressione, oltre alle feature, scaliamo anche
il target $(Z_\text{real}, Z_\text{imag})$: i modelli lineari (Ridge,
Bagging Ridge) ne traggono beneficio in stabilità numerica, e per gli
altri il costo è trascurabile.

\section{Strategia di cross-validation}

Il punto centrale è il concetto di \emph{leakage}: una CV mal progettata
sovrastima sistematicamente le prestazioni perché il modello vede in
addestramento informazioni che non avrebbe a disposizione in produzione.
La nostra regola d'oro è: \emph{lo splitter di CV deve imitare la stessa
procedura di valutazione fuori-sample del task}.

\begin{table}[H]
  \centering
  \caption{Splitter scelti per ciascun task.}
  \label{tab:cv-strategy}
  \begin{tabularx}{\linewidth}{lll X}
    \toprule
    \textbf{Task} & \textbf{Splitter} & \textbf{Gruppo} & \textbf{Motivazione} \\
    \midrule
    1 & \acrshort{logo} & $(Aging, Temp)$ & Il test è una coppia (Aging, Temp) esclusa: la CV deve fare lo stesso. \\
    2 & Hold-out delle 8 curve di una $(Aging, SOC)$ & $(Aging, SOC)$ & Per la coppia scelta dall'utente, le 8 curve di Nyquist (una per temperatura, ${\sim}49$ frequenze ciascuna) vanno interamente in test e sono tutte della stessa classe; la CV interna su training usa StratifiedGroupKFold su $(Aging, Temp)$ per il tuning. \\
    3 & \acrshort{logo} & $Aging$ & Il test è un \emph{intero} livello di Aging escluso: ogni fold pretende di non averlo mai visto. \\
    \bottomrule
  \end{tabularx}
\end{table}

\section{Metriche}

Le metriche sono incapsulate nei dataclass di
\codefile{src/evaluation/metrics.py}:

\begin{itemize}
  \item \textbf{Regressione (Task 1, 3):} \acrshort{mse}, \acrshort{mae},
        $R^{2}$ aggregato e per-componente $R_\text{real}^{2}$,
        $R_\text{imag}^{2}$. La metrica di selezione è
        \acrshort{mse} (Task 1) o $R^{2}$ (Task 3) -- vincono entrambe i
        modelli che meglio bilanciano semicerchio e coda.
  \item \textbf{Classificazione (Task 2):} Accuracy globale sulle 8
        curve di Nyquist (una per temperatura) della coppia
        $(Aging, SOC)$ esclusa, e Accuracy per-temperatura. AUC--ROC
        non è applicabile perché tutte le 8 curve appartengono alla
        stessa classe (Young oppure Old) per costruzione; la confusion
        matrix collassa a una sola riga e quindi non è informativa.
        La metrica di selezione è Accuracy sul hold-out.
\end{itemize}

\section{Hyperparameter tuning}

Il modulo \codefile{src/models/tuning.py} fa GridSearchCV con lo splitter
del task: i grid sono dichiarativi in \codefile{config/config.yaml} sotto
la chiave \texttt{tuning}. Il risultato è persistito in
\codefile{models/tuning/<task>\_best.json} ed è opzionalmente caricabile
dal registry tramite \texttt{regression\_models(use\_tuned=True,
task=\dots)}. Tutti e tre i task seguono lo stesso pattern: si individua
il vincitore del benchmark di default, si lancia il GridSearchCV solo su
quello, si confrontano i due risultati e si tiene la combinazione
migliore.

\section{Riproducibilità}

\begin{itemize}
  \item \texttt{random\_state=42} è centrale e usato in
        \codefile{registry.py}, \codefile{tuning.py} e nei task.
  \item Il loader del dataset ordina deterministicamente prima di
        consegnare i dati.
  \item Le dipendenze sono versionate (con upper bound) in
        \codefile{requirements.txt}.
  \item Ogni run MLflow registra il SHA git, il SHA-256 del CSV, lo
        snapshot della config YAML e i parametri del modello.
\end{itemize}
